% Compilation : pdflatex main.tex  (deux fois, pour la table des matières)
% Les interprétations sont incluses depuis <modèle>/train/interpretation.tex :
% ce fichier ne les duplique pas, il les organise.

\documentclass[11pt,a4paper]{article}
\usepackage[utf8]{inputenc}
\usepackage[T1]{fontenc}
\usepackage[french]{babel}
\usepackage{booktabs,graphicx,geometry,amsmath,hyperref}
\geometry{margin=2.4cm}
\hypersetup{colorlinks=true,linkcolor=black,urlcolor=blue}

\title{Analyse spectrale du graphe d'attention\\
pour la détection d'injection de prompt\\[4pt]
\large Étude exploratoire sur huit modèles --- BIPIA email-QA, split \texttt{train}}
\author{}
\date{14 septembre 2026}

\begin{document}
\maketitle
\tableofcontents
\newpage

% =====================================================================
\section{Objet de ce document}

Ce document rassemble les interprétations rédigées après chaque run de la série.
Les tableaux de résultats, la configuration et les contrôles chiffrés de chaque
run figurent dans le \texttt{rapport.tex} du dossier correspondant, généré
automatiquement à partir des CSV ; ici ne sont reprises que les interprétations,
qui sont rédigées à la main.

\paragraph{Protocole.} Pour chaque paire, le même email est présenté au modèle
dans deux conditions : bénigne, et augmentée d'une instruction d'attaque insérée
au début, au milieu ou à la fin du corps du message. L'insertion étant faite par
le code, la frontière est connue au caractère près et chaque token est étiqueté
\texttt{benign} ou \texttt{injection}. Le prompt complet est envoyé (template de
chat, prompt système, consignes de tâche, question) ; tout ce qui n'est pas
l'attaque est \texttt{benign}. Quarante paires par modèle, soit 80 observations.

\paragraph{Correction pour sélection.} Tous les $p$ rapportés dans ce document
sont corrigés par simulation sous hypothèse nulle, calquée sur le plan exact du
run : on tire des scores indépendants des étiquettes, on prend le maximum sur le
même nombre de candidats, et on compare. Cette correction est indispensable :
avec $2\,880$ candidats et 80 observations, le maximum d'AUROC sous bruit pur
vaut déjà $0{,}72$ en moyenne. Un chiffre non corrigé n'a aucune valeur ici.

\paragraph{Sur les addenda.} Plusieurs interprétations contiennent des addenda
datés qui révisent un énoncé antérieur. C'est délibéré : le texte d'origine est
conservé tel qu'il a été écrit, avec l'information alors disponible, et la
révision est ajoutée à la suite. Trois énoncés ont ainsi été corrigés par des
runs ultérieurs (sections~\ref{sec:smollm3}, \ref{sec:qwen3},
\ref{sec:qwen25-15b}). La chronologie importe davantage que la cohérence
apparente.

\paragraph{Limite majeure, valable pour tout ce document.} \textbf{Aucun de ces
résultats n'a été confirmé sur le split \texttt{test}.} La correction pour
sélection valide l'énoncé « il existe une tête dont $\lambda_2$ sépare les deux
conditions sur ces 40 paires » ; elle ne dit rien de la généralisation à des
catégories d'attaque jamais vues. Les catégories de \texttt{train} et
\texttt{test} étant disjointes dans BIPIA, seul un run sur \texttt{test}, avec
le candidat fixé d'avance, transformerait ces résultats exploratoires en
résultats.

% =====================================================================
\section{Définition des métriques}
\label{sec:metriques}

Toutes les quantités rapportées dans ce document sont définies ici. Elles se
répartissent en six familles selon ce qu'elles utilisent : le graphe seul, le
graphe et les étiquettes, ou une coupe découverte sans étiquettes.

\subsection{Construction du graphe}

L'attention d'un décodeur est causale, donc triangulaire inférieure. Traitée
telle quelle comme un graphe non orienté, toute quantité spectrale serait vide
de sens. Elle est donc symétrisée :
\[
  W = \tfrac{1}{2}\left(A + A^{\top}\right),
\]
où $A$ est la matrice d'attention d'une couche, moyennée sur les têtes
(voie agrégée) ou prise tête par tête (voie \texttt{-{}-per-head}).

Le Laplacien retenu est le \emph{symétrique normalisé} :
\[
  L = I - D^{-1/2} W D^{-1/2},
  \qquad D = \operatorname{diag}\!\left(\textstyle\sum_j W_{ij}\right).
\]
Son spectre est contenu dans $[0, 2]$ quelle que soit la longueur de la
séquence. C'est ce qui rend comparables un prompt bénin et sa variante
injectée, nécessairement plus longue : sans cette normalisation, toute
différence mesurée serait dominée par l'écart de longueur.

On note $0 = \lambda_1 \le \lambda_2 \le \dots \le \lambda_n$ les valeurs
propres de $L$, et $u_2$ le vecteur propre associé à $\lambda_2$, dit
\emph{vecteur de Fiedler}. La décomposition est dense (toutes les valeurs
propres calculées) : un solveur creux tronquerait le spectre et fausserait les
deux métriques définies comme des rapports sur le spectre entier.

\subsection{Métriques spectrales (graphe seul)}

Ces cinq métriques ne dépendent que des poids d'attention. Elles ne requièrent
ni étiquettes ni états cachés, et sont donc calculables sous accès
« attention seule ». Ce sont elles qui composent le tenseur
\texttt{per\_head\_values.npy}.

\begin{description}
\item[\texttt{fiedler\_value}] $\lambda_2$, dite \emph{connectivité
  algébrique}. Nulle si le graphe est déconnecté, d'autant plus petite qu'il
  est proche de l'être. C'est la métrique qui porte le signal dans cette étude.

\item[\texttt{connectivity\_ratio}] $\lambda_2 / \lambda_{\max}$. Même
  information que $\lambda_2$, rapportée à l'étendue du spectre, donc
  insensible à un facteur d'échelle global. Quasi redondante avec
  \texttt{fiedler\_value} lorsque $\lambda_{\max}$ varie peu, ce qui explique
  que les deux apparaissent systématiquement ensemble dans les classements.

\item[\texttt{spectral\_radius}] $\lambda_{\max}$, la plus grande valeur
  propre.

\item[\texttt{spectral\_entropy\_norm}] entropie de la distribution des valeurs
  propres, normalisée :
  \[
    H = -\frac{1}{\log n}\sum_{k=1}^{n} p_k \log p_k,
    \qquad p_k = \frac{\lambda_k}{\sum_j \lambda_j}.
  \]
  Proche de 1 si le spectre est étalé uniformément, plus faible s'il est
  concentré sur quelques modes.

\item[\texttt{hfer}] \emph{high-frequency energy ratio} : part de la masse
  spectrale portée par la moitié haute du spectre,
  \[
    \mathrm{hfer} = \frac{\sum_{k > n/2} \lambda_k}{\sum_k \lambda_k}.
  \]
  Élevée quand le graphe porte beaucoup de variation à courte échelle.
\end{description}

\noindent Ces cinq noms, leur ordre et leurs définitions sont ceux de
\texttt{spectral\_trust.per\_head.PER\_HEAD\_METRIC\_NAMES}, la bibliothèque
qui les calcule effectivement, de sorte que la voie agrégée et la voie par tête
sont directement comparables.

\subsection{Statistiques de coupe (graphe et étiquettes)}

Ces quantités utilisent la vraie frontière entre tokens bénins ($\mathcal{B}$)
et tokens injectés ($\mathcal{I}$). Elles décrivent la structure, elles ne
constituent pas un détecteur. Dans ce qui suit, $A$ est $W$ privée de sa
diagonale.

\begin{description}
\item[\texttt{n\_benign}, \texttt{n\_injection}] effectifs des deux familles
  après filtrage des n\oe uds.

\item[\texttt{density\_benign}, \texttt{density\_injection}] poids moyen
  \emph{par paire possible} à l'intérieur de chaque famille :
  \[
    d_{\mathcal{B}}
      = \frac{\sum_{i \ne j \in \mathcal{B}} A_{ij}}
             {|\mathcal{B}|\,(|\mathcal{B}|-1)}.
  \]
  La normalisation par le nombre de paires n'est pas cosmétique :
  l'injection est presque toujours beaucoup plus courte que l'hôte
  (ici $4\,\%$ des tokens), et une somme brute favoriserait mécaniquement le
  groupe le plus grand.

\item[\texttt{density\_cross}] même définition entre les deux familles,
  $d_{\times} = \left(\sum_{i \in \mathcal{B}, j \in \mathcal{I}} A_{ij}\right)
  / (|\mathcal{B}|\,|\mathcal{I}|)$.

\item[\texttt{separation}] l'hypothèse du bicluster en un nombre :
  \[
    \mathrm{sep} = \frac{\sqrt{d_{\mathcal{B}} \cdot d_{\mathcal{I}}}}{d_\times}.
  \]
  Supérieure à 1, les deux familles se tiennent davantage entre elles qu'elles
  ne se parlent. \textbf{Attention} : cette quantité utilise les étiquettes.
  Elle constate une structure, elle ne mesure pas une structure
  \emph{exploitable} sans étiquettes --- la section~\ref{sec:synthese} montre
  qu'elle ne prédit pas où un détecteur fonctionne.

\item[\texttt{cut\_weight}] poids total de la coupe entre les deux familles,
  $\mathrm{cut} = 2\sum_{i \in \mathcal{B}, j \in \mathcal{I}} A_{ij}$
  (le facteur 2 compte l'arête dans les deux sens, le graphe étant non
  orienté).

\item[\texttt{conductance}] $\mathrm{cut} / \min(\mathrm{vol}_{\mathcal{B}},
  \mathrm{vol}_{\mathcal{I}})$, où $\mathrm{vol}$ est la somme des degrés du
  groupe. Petite quand la coupe est nette.

\item[\texttt{normalized\_cut}]
  $\mathrm{cut}/\mathrm{vol}_{\mathcal{B}} + \mathrm{cut}/\mathrm{vol}_{\mathcal{I}}$,
  variante qui pénalise les groupes déséquilibrés.

\item[\texttt{modularity}] modularité de Newman de la bipartition
  bénin/injection :
  \[
    Q = \sum_{c \in \{\mathcal{B},\mathcal{I}\}}
        \left[\frac{e_c}{2m} - \left(\frac{d_c}{2m}\right)^{2}\right],
  \]
  où $e_c$ est le poids interne au groupe $c$, $d_c$ la somme de ses degrés et
  $2m$ le poids total. Positive si les groupes sont plus denses qu'un graphe
  aléatoire de mêmes degrés.
\end{description}

\subsection{Coupe découverte (sans étiquettes)}

Un détecteur ne sait pas où se trouve l'injection. Il découvre une coupe et en
mesure la qualité. La coupe retenue ici est donnée par le \emph{signe du
vecteur de Fiedler} : $\mathcal{S} = \{i : u_2(i) \ge 0\}$.

\begin{description}
\item[\texttt{cut\_conductance}] conductance de cette coupe découverte. C'est
  le score du détecteur, et il est calculable aussi sur un prompt bénin, où la
  meilleure coupe disponible reste mauvaise --- c'est précisément là que
  devrait résider le signal.

\item[\texttt{cut\_modularity}] modularité de Newman de la même coupe
  découverte.
\end{description}

\noindent \textbf{Limite structurelle.} La bipartition par le signe de $u_2$
produit deux groupes de taille comparable. Elle ne peut donc pas isoler une
minorité de $4\,\%$, ce qui explique l'échec de \texttt{cut\_conductance}
rapporté en section~\ref{sec:synthese}. Une \emph{coupe par balayage} --- trier
les n\oe uds selon $u_2$ et retenir le seuil minimisant la conductance ---
lèverait cette limite ; elle n'est pas implémentée dans cette série.

\subsection{Accord entre la coupe découverte et la vérité}

\begin{description}
\item[\texttt{accuracy}] proportion de n\oe uds correctement classés, prise au
  maximum sur les deux orientations d'étiquettes, le signe d'un vecteur propre
  étant arbitraire.

  \textbf{Cette métrique se lit mal et doit toujours être accompagnée de sa
  ligne de base.} Parce qu'elle maximise sur les deux orientations, placer tous
  les n\oe uds dans un seul groupe marque déjà $\max(p, 1-p)$, où $p$ est la
  fraction injectée. Avec $p \approx 0{,}04$, cette ligne de base vaut
  $\approx 0{,}96$ : un accord observé de $0{,}58$ est donc très
  \emph{inférieur} au trivial, et non « un peu au-dessus du hasard ».

\item[\texttt{ari}] indice de Rand ajusté entre la coupe découverte et la vraie
  partition. Corrigé du hasard par construction : $0$ correspond à une coupe
  sans rapport avec la frontière, $1$ à une correspondance parfaite, une valeur
  négative à un accord inférieur au hasard. \textbf{C'est la métrique d'accord
  à rapporter}, et non \texttt{accuracy}.
\end{description}

\subsection{Agrégation sur l'ensemble des paires}

\begin{description}
\item[\texttt{auroc}] aire sous la courbe ROC, calculée par les rangs
  (statistique de Mann--Whitney), les ex \ae quo recevant le rang moyen. Elle
  répond à : « en classant les 80 observations selon ce score, les prompts
  injectés remontent-ils ? ». $0{,}5$ correspond à l'absence de séparation.

\item[\texttt{direction}] $+1$ si un score élevé indique l'injection, $-1$
  sinon. L'AUROC rapportée est corrigée en conséquence, de sorte qu'elle est
  toujours $\ge 0{,}5$ ; c'est pourquoi toute valeur doit être jugée par
  rapport au maximum sous hypothèse nulle et non à $0{,}5$.
\end{description}

\subsection{Quantités descriptives}

\begin{description}
\item[\texttt{n\_nodes}] n\oe uds conservés dans le graphe après filtrage.
\item[\texttt{n\_tokens}] longueur du prompt en tokens.
\item[\texttt{bandwidth}] distance moyenne $|i-j|$ pondérée par les poids
  d'attention. Faible quand l'attention est locale (une bande autour de la
  diagonale), élevée quand elle porte des liens à longue portée. Sert à
  caractériser le régime d'une couche, non à détecter.
\end{description}

% =====================================================================
\section{Synthèse sur les huit modèles}
\label{sec:synthese}

\begin{table}[htbp]\centering\small
\begin{tabular}{lrrlrrlr}
\hline
modèle & L & H & meilleur triplet par tête & AUROC & $p$ & agrégé $\lambda_2$ & $p$ \\
\hline
Llama-3.2-1B  & 16 & 32 & L6, t.21, \texttt{conn.}   & 0,884 & $<$0,004 & L0 \ 0,604 & 0,97 \\
Qwen2.5-1.5B  & 28 & 12 & L9, t.9, \texttt{conn.}    & 0,813 & $<$0,004 & L19 0,757 & \textbf{0,003} \\
SmolLM3-3B    & 36 & 16 & L24, t.0, \texttt{radius}  & 0,791 & 0,007    & L0 \ 0,759 & \textbf{0,000} \\
Llama-3.2-3B  & 28 & 24 & L14, t.4, \texttt{fiedler} & \textbf{0,917} & $<$0,004 & L4 \ 0,676 & 0,22 \\
Qwen3-1.7B    & 28 & 16 & L17, t.6, \texttt{fiedler} & 0,907 & $<$0,004 & L20 0,657 & 0,55 \\
SmolLM2-1.7B  & 24 & 32 & L17, t.3, \texttt{fiedler} & 0,858 & $<$0,004 & L0 \ 0,624 & 0,95 \\
gemma-2-2b-it & 26 &  8 & L9, t.0, \texttt{conn.}    & 0,779 & $<$0,004 & L0 \ 0,611 & 0,99 \\
Qwen2.5-3B    & 36 & 16 & L18, t.10, \texttt{fiedler}& 0,876 & $<$0,004 & L9 \ 0,798 & \textbf{0,000} \\
\hline
\end{tabular}
\caption{Résultats par modèle, dans l'ordre chronologique des runs.
\texttt{conn.} $=$ \texttt{connectivity\_ratio}, \texttt{fiedler} $=$
\texttt{fiedler\_value}, \texttt{radius} $=$ \texttt{spectral\_radius}.
L $=$ couches, H $=$ têtes.}
\label{tab:synthese}
\end{table}

\subsection{Ce qui est établi}

\paragraph{1. \texttt{cut\_conductance} ne détecte rien : 8 modèles sur 8.}
Les huit meilleures AUROC tiennent dans $[0{,}613\,;\,0{,}697]$, la plage même
du maximum sous hypothèse nulle. Le plus petit $p$ corrigé vaut $0{,}04$ (un
seul modèle), tous les autres dépassent $0{,}18$.

C'est une réfutation de la prémisse initiale du travail. Le « détecteur sans
label » fondé sur la qualité de la coupe découverte par le signe du vecteur de
Fiedler ne fonctionne pas sur les injections indirectes de BIPIA. Ce résultat
négatif est solide et mérite d'être publié comme tel.

\paragraph{2. Le mécanisme de cet échec est identifié.}
L'indice de Rand ajusté entre la coupe découverte et la vraie frontière vaut le
hasard sur les huit modèles ($-0{,}011$ à $+0{,}003$). La cause est
structurelle : la bipartition par le signe de $u_2$ produit deux groupes de
taille comparable, alors que l'injection représente $3{,}6$ à $4{,}1\,\%$ des
tokens. Une coupe équilibrée ne peut pas isoler une telle minorité.

L'échec n'est donc pas une absence de signal dans le graphe, mais une
inadéquation de la règle de coupe à la cible. La modification indiquée est le
passage à une \emph{coupe par balayage} --- trier les n\oe uds selon $u_2$ et
retenir le seuil qui minimise la conductance --- qui autorise des coupes
déséquilibrées et reste bornée par l'inégalité de Cheeger.

\paragraph{3. La connectivité algébrique porte le signal : 8 modèles sur 8.}
Le meilleur triplet par tête est significatif après correction sur les huit
modèles, avec des AUROC de $0{,}779$ à $0{,}917$ (médiane $0{,}867$). Sept des
huit désignent une métrique de la famille $\lambda_2$
(\texttt{fiedler\_value} $=\lambda_2$ ou
\texttt{connectivity\_ratio} $=\lambda_2/\lambda_{\max}$) ; seul SmolLM3-3B
retient \texttt{spectral\_radius}.

Sur six modèles, la structure du classement renforce ce constat : deux ou trois
têtes distinctes se détachent, chacune confirmée indépendamment par les deux
métriques $\lambda_2$, et la même tête réapparaît parfois à plusieurs
profondeurs. Un maximum isolé peut être un accident de sélection ; cette
structure ne l'est pas.

\paragraph{4. Les deux voies convergent.}
Trois modèles présentent aussi un signal \emph{agrégé} significatif, toujours
via \texttt{fiedler\_value}, jamais via \texttt{cut\_conductance}. Sur
Qwen2.5-3B, trois des quatre scores agrégés culminent à la même couche L9.

Le fait que la voie par tête et la voie agrégée désignent la même quantité,
$\lambda_2$, alors qu'elles reposent sur des calculs distincts, est l'argument
le plus difficile à attribuer au hasard. La conclusion à retenir est que
$\lambda_2$ est le descripteur pertinent, et que \texttt{cut\_conductance},
autour duquel le projet avait été conçu, est la mauvaise statistique.

\subsection{Ce qui n'est pas établi}

\paragraph{La localisation ne se réplique pas.} Ni la couche, ni la tête ne se
retrouvent d'un modèle à l'autre. Il n'existe pas de « tête à injection »
universelle : chaque modèle doit être sondé séparément. Seule la famille de
métriques se réplique.

\paragraph{La part de contiguïté n'est pas quantifiée.} La séparation décroît
avec la profondeur sur sept modèles sur huit ($r$ de $-0{,}24$ à $-0{,}75$),
avec un contre-exemple franc (SmolLM2-1.7B, $r = +0{,}46$). Cette tendance
suggère que \texttt{separation} mesure largement la contiguïté des tokens
injectés plutôt que leur nature --- si elle mesurait un groupement sémantique,
on l'attendrait croissante avec la profondeur. Le contrôle direct, consistant à
insérer dans un prompt bénin un segment factice de même longueur et à mesurer
la séparation obtenue, n'est pas implémenté. Tant qu'il manque, l'argument
reste indirect.

\paragraph{\texttt{separation} ne prédit pas la détection.} La corrélation entre
séparation moyenne et AUROC de détection, couche par couche, va de $-0{,}30$ à
$+0{,}50$ selon le modèle, sans régularité. \texttt{separation} utilise les
vraies étiquettes et constate une structure ; elle ne mesure pas une structure
exploitable sans étiquettes. Les deux quantités ne doivent pas être présentées
comme deux facettes du même phénomène.

\paragraph{Le nombre de têtes : piste non concluante.} La marge du meilleur
triplet au-dessus du seuil de bruit corrèle à $r = +0{,}55$ avec le nombre de
têtes du modèle. Avec $n = 8$, ce n'est pas significatif, et les deux modèles à
16 têtes affichent les marges les plus écartées du tableau. À mentionner comme
observation, pas comme résultat.

\subsection{Biais connus du protocole}

Le puits d'attention (token BOS) capte une part importante de la masse et se
trouve, dans ce protocole, étiqueté \texttt{benign} : il gonfle
\texttt{density\_benign} sur les huit modèles. Le panneau circulaire des figures
place les n\oe uds en ordre de lecture, de sorte qu'une injection contiguë
occupe un arc contigu par construction ; seules la matrice d'adjacence
réordonnée et les statistiques de coupe constituent une preuve. Enfin, l'accord
brut (\texttt{accuracy}) doit toujours être comparé à sa ligne de base triviale,
de l'ordre de $0{,}96$ ici : l'indice de Rand ajusté est la métrique à
rapporter.

\subsection{Étape suivante}

Huit hypothèses sont pré-enregistrées, une par modèle : le triplet
(couche, tête, métrique) de la table~\ref{tab:synthese}. Elles doivent être
rapportées sur le split \texttt{test}, dont les quinze catégories d'attaque sont
disjointes de celles de \texttt{train}, en lisant \emph{uniquement} la ligne
correspondante et sans consulter le maximum du nouveau run. Un seul candidat par
modèle, donc aucune correction à appliquer.

Le candidat agrégé de Qwen2.5-3B (\texttt{fiedler\_value} à L9, $0{,}798$)
mérite d'être testé en priorité : un détecteur agrégé n'exige ni l'identification
d'une tête, ni le stockage d'un tenseur de \emph{features}, et serait donc plus
simple à défendre et à déployer que la voie par tête.

% =====================================================================
\newpage
\section{Interprétations par modèle}

Les sections qui suivent reproduisent, sans modification, les interprétations
rédigées après chaque run, dans l'ordre chronologique.

\subsection{Llama-3.2-1B-Instruct}
\label{sec:llama1b}
\InputIfFileExists{Llama-3.2-1B-Instruct/train/interpretation.tex}{}{%
\textit{Fichier d'interprétation absent.}}

\subsection{Qwen2.5-1.5B-Instruct}
\label{sec:qwen25-15b}
\InputIfFileExists{Qwen2.5-1.5B-Instruct/train/interpretation.tex}{}{%
\textit{Fichier d'interprétation absent.}}

\subsection{SmolLM3-3B}
\label{sec:smollm3}
\InputIfFileExists{SmolLM3-3B/train/interpretation.tex}{}{%
\textit{Fichier d'interprétation absent.}}

\subsection{Llama-3.2-3B-Instruct}
\label{sec:llama3b}
\InputIfFileExists{Llama-3.2-3B-Instruct/train/interpretation.tex}{}{%
\textit{Fichier d'interprétation absent.}}

\subsection{Qwen3-1.7B}
\label{sec:qwen3}
\InputIfFileExists{Qwen3-1.7B/train/interpretation.tex}{}{%
\textit{Fichier d'interprétation absent.}}

\subsection{SmolLM2-1.7B-Instruct}
\label{sec:smollm2}
\InputIfFileExists{SmolLM2-1.7B-Instruct/train/interpretation.tex}{}{%
\textit{Fichier d'interprétation absent.}}

\subsection{gemma-2-2b-it}
\label{sec:gemma}
\InputIfFileExists{gemma-2-2b-it/train/interpretation.tex}{}{%
\textit{Fichier d'interprétation absent.}}

\subsection{Qwen2.5-3B-Instruct}
\label{sec:qwen25-3b}
\InputIfFileExists{Qwen2.5-3B-Instruct/train/interpretation.tex}{}{%
\textit{Fichier d'interprétation absent.}}

\end{document}
